% =====================================================================
\chapter{Teorema fundamental del cálculo, áreas y valor medio}
\label{tema:2}
% =====================================================================

El \cref{tema:1} definió la integral como un límite de sumas, una definición rigurosa pero poco práctica: calcular $\int_0^1x^{3}\dif x$ por sumas ya exige conocer $\sum k^{3}$. El teorema fundamental del cálculo cambia por completo la situación, porque reduce el cálculo de una integral al de una primitiva y, a la vez, garantiza que toda función continua tiene primitiva. Este tema lo demuestra en sus dos partes y lo aplica a áreas, valores medios y funciones definidas por integrales, como el logaritmo. Se siguen el capítulo~14 de Spivak~\cite{spivak}, el capítulo~5 de Apostol~\cite{apostol}, la §5.3 de~\cite{openstax1}, las §5.6 y~5.7 de Strang~\cite{strang} y las partes~A a~C de la unidad~3 del curso 18.01SC~\cite{mit1801}.

\section{El teorema del valor medio integral}

\begin{teorema}[Valor medio integral]\label{teo:tvm-integral}
Sea $f$ continua en $[a,b]$. Existe $c\in[a,b]$ tal que
\[
\int_a^b f(x)\dif x=f(c)\,(b-a).
\]
\end{teorema}

\begin{definicion}[Valor medio de una función]\label{def:valor-medio}
Sea $f$ integrable en $[a,b]$. Su \emph{valor medio} en $[a,b]$ es
$\displaystyle\bar f=\frac{1}{b-a}\int_a^b f(x)\dif x$.
\end{definicion}

Geométricamente, si $f\ge0$, el rectángulo de base $[a,b]$ y altura $\bar f$ tiene la misma área que la región bajo la gráfica. El \cref{teo:tvm-integral} dice que una función continua toma su valor medio en algún punto. El nombre procede de la analogía con la media de $n$ números: $\bar f$ es el límite de las medias $\frac1n\sum_{k=1}^{n}f(x_k)$ en puntos equiespaciados, que son sumas de Riemann divididas por $b-a$.

\section{El teorema fundamental del cálculo}

\begin{teorema}[Teorema fundamental del cálculo, primera parte]\label{teo:tfc1}
Sea $f$ integrable en $[a,b]$ y sea $F(x)=\int_a^x f(t)\dif t$ para $x\in[a,b]$.
\begin{enumerate}
  \item $F$ es continua en $[a,b]$; de hecho, si $\abs{f}\le K$, entonces $\abs{F(y)-F(x)}\le K\abs{y-x}$.
  \item Si $f$ es continua en un punto $c\in(a,b)$, entonces $F$ es derivable en $c$ y $F'(c)=f(c)$.
\end{enumerate}
\end{teorema}

\begin{corolario}\label{cor:existencia-primitiva}
Toda función continua en un intervalo $I$ tiene primitiva en $I$: para cualquier $a\in I$, $F(x)=\int_a^x f(t)\dif t$ lo es.
\end{corolario}

\begin{corolario}[Derivación de integrales con límites variables]\label{cor:leibniz}
Sea $f$ continua en un intervalo $I$ y sean $g,h$ derivables con valores en $I$. Entonces
\[
\frac{d}{dx}\int_{h(x)}^{g(x)}f(t)\dif t=f\bigl(g(x)\bigr)\,g'(x)-f\bigl(h(x)\bigr)\,h'(x).
\]
\end{corolario}

\begin{teorema}[Teorema fundamental del cálculo, segunda parte, o regla de Barrow]\label{teo:barrow}
Sea $f$ integrable en $[a,b]$ y sea $F$ continua en $[a,b]$, derivable en $(a,b)$, con $F'(x)=f(x)$ para todo $x\in(a,b)$. Entonces
\[
\int_a^b f(x)\dif x=F(b)-F(a)=\eval{F(x)}{a}{b}.
\]
\end{teorema}

La demostración no exige que $f$ sea continua, solo integrable, y no usa la primera parte. Cuando $f$ es continua también puede deducirse de ella: $G(x)=\int_a^x f$ y $F$ son primitivas de $f$ en $(a,b)$, continuas en $[a,b]$, luego difieren en una constante (\cref{prop:primitivas-constante}), y $F(b)-F(a)=G(b)-G(a)=\int_a^b f$.

\addtocounter{definicion}{1}%

\section{El logaritmo como integral}

El \cref{cor:existencia-primitiva} permite \emph{definir} funciones mediante integrales. El ejemplo más importante es el logaritmo, que puede construirse sin conocer previamente la exponencial~\cite[cap.~18]{spivak}, \cite[cap.~6]{apostol}.

\begin{proposicion}\label{prop:logaritmo}
Sea $L(x)=\int_1^x\frac{\dif t}{t}$ para $x>0$. Entonces $L(1)=0$, $L'(x)=\frac1x$, $L$ es estrictamente creciente y $L(ab)=L(a)+L(b)$ para todos $a,b>0$.
\end{proposicion}

$L$ es, por tanto, el logaritmo natural, $L=\ln$, y su inversa es la exponencial. De la misma forma se definen funciones que no tienen primitiva expresable con funciones elementales, como la función de error $\operatorname{erf}(x)=\frac{2}{\sqrt\pi}\int_0^x e^{-t^{2}}\dif t$ o la función de distribución de la normal estándar, $\Phi(x)=\frac{1}{\sqrt{2\pi}}\int_{-\infty}^{x}e^{-t^{2}/2}\dif t$, que es una integral impropia (\cref{tema:4}). Por el \cref{teo:tfc1}, $\Phi'(x)=\frac{1}{\sqrt{2\pi}}e^{-x^{2}/2}$ aunque $\Phi$ no tenga fórmula elemental.

\section{Áreas entre curvas}

\begin{definicion}[Área entre dos curvas]\label{def:area}
Sean $f,g$ continuas en $[a,b]$. El área de la región comprendida entre sus gráficas y las rectas $x=a$, $x=b$ es $\int_a^b\abs{f(x)-g(x)}\dif x$. Si $f\ge g$ en $[a,b]$, es $\int_a^b(f-g)$.
\end{definicion}

Para calcularla se hallan los puntos donde $f-g$ cambia de signo, se divide $[a,b]$ en subintervalos donde el signo es constante y se suman las integrales de $f-g$ o de $g-f$ según corresponda.

\addtocounter{definicion}{1}%

\addtocounter{definicion}{1}%

\section{Valor medio y simetrías}

\addtocounter{definicion}{1}%

\begin{proposicion}[Integrales de funciones pares e impares]\label{prop:par-impar}
Sea $f$ continua en $[-a,a]$. Si $f$ es par, $\int_{-a}^{a}f=2\int_0^{a}f$; si $f$ es impar, $\int_{-a}^{a}f=0$.
\end{proposicion}

\begin{observacion}[Medias ponderadas]
Si $w\ge0$ es continua con $\int_a^b w=1$, el número $\int_a^b f(x)\,w(x)\dif x$ es una media de $f$ ponderada por $w$. Cuando $w$ es la densidad de una variable aleatoria $X$, es la esperanza $\E[f(X)]$ del módulo~5; el valor medio de la \cref{def:valor-medio} corresponde a la densidad uniforme $w=\frac{1}{b-a}$. El \cref{eje:2-tvm-ponderado} generaliza el \cref{teo:tvm-integral} a este caso.
\end{observacion}

\addtocounter{definicion}{1}%

\addtocounter{definicion}{1}%

